\section{结论}\label{sec:conclusion}

本文构造了一种基于通行的一维目标相对历史压缩，并在所得账户空间上建立可加测量理论。通用载体识别可加对象，精确纤维定理则刻画保持该对象不变的可行时间顺序变换。在有限步一维定义域中，有序端点与总跨越剖面随后可以重构完整有向账本。

BPCC 把实际历史比较转化为精确的偏可加次序。PVBC 刻画何时阿基米德完备化保持实际历史次序，从而允许用归一化实数支撑赋值的一致同意作精确表示。有限网格的字典序例子说明：即使物理状态网格有限，只要路径可无界重复，这一限制仍具有实质内容。PROC 提供一个对边界敏感的正则性细化，从而得到允许目标原子的有向 Stieltjes 测度。

测度表示逐个赋值给出“端点--总通行--目标达成”的解释。经济学内容则在原始排名与识别集合层面表述。对于任意有限个具有共同端点、且在应用上可行的遍历，只要匹配时间安排足够长，就可以仅通过增加端点停留，在不改变遍历速度的情况下固定持续时间和一个能够区分端点的连续累计流量统计量。相同账户差分可以跨环境转移序数排名，有限个已观察选择则对反事实通行对比给出锐界；方案--价格数据还可以把同一界表示为货币单位。

通行账户有意把可交换的闭合游程视为在时间顺序上中性。纤维定理使这一限制精确化：它识别压缩究竟丢失了哪些信息，而不是把所有终点相同的历史都视为等价。若要评价闭合游程的顺序、首次达成效应、内生目标或动态状态转移，则需要更丰富的记录；这些也构成超越本文可加通行域的自然扩展。
